% Reconstructed from the supplied figure; original raster icons are in assets/.
% Compile with pdfLaTeX: pdflatex -interaction=nonstopmode -halt-on-error figure_biolinum.tex
% This Linux Biolinum version is Figure 1 of paper/AAMAS/main.tex (fig:overview).
% All headings, labels, numbers, and mathematical symbols are real PDF text.
% Canvas: 188 x 90 mm (compressed 19-09-2026 from the 188 x 106 mm layout of figure.tex so the
% paper keeps its 8-page body; text sizes are unchanged except the two panel titles).
% Coordinates below are millimetres, y increases downward.
\documentclass[tikz,border=0pt]{standalone}
\usepackage[T1]{fontenc}
\usepackage[utf8]{inputenc}
\usepackage{lmodern}
\usepackage{biolinum}
\renewcommand{\familydefault}{\sfdefault}
\usepackage{amsmath}
\usepackage{graphicx}
\usepackage{tikz}
\usetikzlibrary{calc}
\input{glyphtounicode}
\pdfgentounicode=1
\pdfcompresslevel=9
\pdfobjcompresslevel=2
\pdfinfo{/Title (Collective-risk game and four experimental designs)
         /Subject (TikZ reconstruction with selectable text and extracted transparent icons)
         /Creator (pdfLaTeX and TikZ)}

% Flat, print-friendly colours. Edit these definitions to restyle the figure.
\definecolor{Ink}{HTML}{102033}
\definecolor{Muted}{HTML}{415166}
\definecolor{Blue}{HTML}{0867B3}
\definecolor{Green}{HTML}{008B64}
\definecolor{BlueLine}{HTML}{38AFE0}
\definecolor{GreenLine}{HTML}{24B887}
\definecolor{PanelWash}{HTML}{F3FAFE}
\definecolor{ModelWash}{HTML}{F5FAFD}
\definecolor{ClaudeWash}{HTML}{FFF0DF}
\definecolor{GPTWash}{HTML}{E5F8ED}
\definecolor{GeminiWash}{HTML}{E6F4FE}
\definecolor{QwenWash}{HTML}{F0E8FB}
\definecolor{GrokWash}{HTML}{FFF7DB}
\definecolor{ClaudeLine}{HTML}{E5AE75}
\definecolor{GPTLine}{HTML}{70CA99}
\definecolor{GeminiLine}{HTML}{79BFEA}
\definecolor{QwenLine}{HTML}{BC94E4}
\definecolor{GrokLine}{HTML}{DDC065}
\definecolor{SelfWash}{HTML}{E7F5FD}
\definecolor{PromptWash}{HTML}{FFF0E1}
\definecolor{ScriptWash}{HTML}{F0E7FB}
\definecolor{MixedWash}{HTML}{FFF7DE}
\definecolor{Orange}{HTML}{C5510A}
\definecolor{Purple}{HTML}{7832B1}
\definecolor{Gold}{HTML}{B07808}
\definecolor{SuccessWash}{HTML}{DCF8E8}
\definecolor{FailureWash}{HTML}{FFE6E4}
\definecolor{SurvivalWash}{HTML}{FFF5D7}
\definecolor{SuccessLine}{HTML}{58C588}
\definecolor{FailureLine}{HTML}{EA9691}
\definecolor{SurvivalLine}{HTML}{E2C56B}
\definecolor{FixedFill}{HTML}{65706F}

% Colour variants: palettes/<name>.tex may redefine any colour above or any
% role below; build them all with palettes/build_palettes.sh. Without
% \FigPalette this file compiles to the paper's Figure 1 unchanged.
\newif\ifheadrule \headrulefalse
\ifdefined\FigPalette \input{\FigPalette}\fi
\makeatletter
\newcommand{\defaultcolorlet}[2]{\@ifundefinedcolor{#1}{\colorlet{#1}{#2}}{}}
\makeatother
\defaultcolorlet{PanelFill}{white}
\defaultcolorlet{HeadA}{Blue}
\defaultcolorlet{HeadB}{Green}
\defaultcolorlet{HeadText}{white}
\defaultcolorlet{HeadTagA}{HeadText}
\defaultcolorlet{HeadTagB}{HeadText}
\defaultcolorlet{HeadRuleA}{BlueLine}
\defaultcolorlet{HeadRuleB}{GreenLine}
\defaultcolorlet{CardLine}{GeminiLine}
\defaultcolorlet{RuleCard}{white}
\defaultcolorlet{SelfLine}{GeminiLine}
\defaultcolorlet{PromptLine}{ClaudeLine}
\defaultcolorlet{ScriptLine}{QwenLine}
\defaultcolorlet{MixedLine}{GrokLine}
\defaultcolorlet{NumOne}{Blue}
\defaultcolorlet{NumTwo}{Orange}
\defaultcolorlet{NumThree}{Purple}
\defaultcolorlet{NumFour}{Gold}
\defaultcolorlet{NumText}{white}
\defaultcolorlet{NumTextOne}{NumText}
\defaultcolorlet{NumTextTwo}{NumText}
\defaultcolorlet{NumTextThree}{NumText}
\defaultcolorlet{NumTextFour}{NumText}
\defaultcolorlet{FixedText}{white}

\newcommand{\fsize}[1]{\fontsize{#1}{\dimexpr #1pt*6/5\relax}\selectfont}
\newcommand{\card}[6]{%
  \path[rounded corners=1.25mm,fill=#5,draw=#6,line width=.36pt]
    (#1,#2) rectangle (#3,#4);}
% #1 filename stem, #2 x centre, #3 y centre, #4 max width, #5 max height.
\newcommand{\asset}[5]{%
  \node[inner sep=0pt,outer sep=0pt] at (#2,#3)
    {\includegraphics[width=#4mm,height=#5mm,keepaspectratio]{assets/#1.png}};}
% A small typographic badge; no raster text is used for these five labels.
\newcommand{\fixedbadge}[2]{%
  \node[rounded corners=.65mm,fill=FixedFill,text=FixedText,inner xsep=.65mm,
        inner ysep=.42mm,outer sep=0pt,font=\fsize{5.8}]
        at (#1,#2) {Fixed};}
% Header block for an experimental-design card; #8 is the number's colour.
\newcommand{\designhead}[8]{%
  \path[fill=#4] (#1,#2) circle[radius=3.8mm];
  \node[text=#8,font=\fsize{14}\bfseries,inner sep=0pt] at (#1,#2) {#3};
  \node[text=Ink,font=\fsize{8.5}\bfseries,inner sep=0pt,align=center]
    at (#5,{#2-1.50}) {#6};
  \node[text=Muted,font=\fsize{7.4},inner sep=0pt,align=center]
    at (#5,{#2+1.80}) {#7};}

\begin{document}
\begin{tikzpicture}[x=1mm,y=-1mm,text=Ink,
                    every node/.style={outer sep=0pt}]
  \path[use as bounding box] (0,0) rectangle (188,90);

  % -------------------- (a) THE GAME --------------------
  \path[rounded corners=1.8mm,fill=PanelFill] (1,1) rectangle (92,89);
  \begin{scope}
    \clip[rounded corners=1.8mm] (1,1) rectangle (92,89);
    \fill[HeadA] (1,1) rectangle (92,8.8);
  \end{scope}
  \ifheadrule \draw[HeadRuleA,line width=.55pt] (1,8.8) -- (92,8.8); \fi
  \path[rounded corners=1.8mm,draw=BlueLine,line width=.55pt]
    (1,1) rectangle (92,89);
  \node[anchor=west,text=HeadText,font=\fsize{14}\bfseries,inner sep=0pt]
    at (3.5,4.9) {\textcolor{HeadTagA}{(a)}\quad The game};

  \card{3}{10.6}{55.5}{60.4}{PanelWash}{CardLine}
  \card{57.5}{10.6}{90}{60.4}{RuleCard}{CardLine}
  \asset{group_game}{29.25}{35.5}{45}{45}

  % Six rule icons and their selectable labels (pitch 7.8 mm).
  \asset{icon_players}{62.8}{15.9}{7.4}{6.7}
  \asset{icon_rounds}{62.8}{23.7}{6.5}{6.7}
  \asset{icon_endowment}{62.8}{31.5}{6.5}{6.7}
  \asset{icon_contribution}{62.8}{39.3}{6.5}{6.7}
  \asset{icon_target}{62.8}{47.1}{6.5}{6.7}
  \asset{icon_risk}{62.8}{54.9}{6.5}{6.9}
  \foreach \y/\label in {15.9/{6 players},23.7/{10 rounds},31.5/{40 each},
                          39.3/{pay 0 / 2 / 4},47.1/{target 120}} {
    \node[anchor=west,inner sep=0pt,font=\fsize{8.4}]
      at (68.2,\y) {\label};
  }
  \node[anchor=west,inner sep=0pt,font=\fsize{7.4}]
    at (68.2,53.7) {miss $\to$ lose all};
  \node[anchor=west,inner sep=0pt,text=Muted,font=\fsize{6.8}]
    at (68.2,56.9) {(probability $p$)};

  % Outcomes: corrected parent heading. The original put 'Target reached'
  % under 'if total < 120'. Here 'miss' has its own explicit condition.
  \card{3}{62}{90}{87}{PanelWash}{CardLine}
  \node[anchor=west,font=\fsize{9.2}\bfseries,inner sep=0pt]
    at (4.4,65.0) {Outcomes};
  \node[anchor=east,font=\fsize{6.8},text=Muted,inner sep=0pt]
    at (88.7,65.0) {Miss: total $<120$};

  \card{4}{67.3}{31.667}{86}{SuccessWash}{SuccessLine}
  \card{32.667}{67.3}{60.333}{86}{FailureWash}{FailureLine}
  \card{61.333}{67.3}{89}{86}{SurvivalWash}{SurvivalLine}
  \asset{icon_success}{17.8335}{71.4}{6.2}{6.2}
  \asset{icon_catastrophe}{46.5}{71.4}{6.2}{6.2}
  \asset{icon_no_catastrophe}{75.1665}{71.4}{6.2}{6.2}

  \node[font=\fsize{7.4}\bfseries,inner sep=0pt] at (17.8335,76.9) {Target reached};
  \node[font=\fsize{7.0}\bfseries,inner sep=0pt] at (46.5,76.9) {Miss + catastrophe};
  \node[font=\fsize{7.0}\bfseries,inner sep=0pt] at (75.1665,76.9) {Miss + no catastrophe};
  \node[font=\fsize{6.8},text=Muted,inner sep=0pt] at (17.8335,80.3) {total $\geq120$};
  \node[font=\fsize{6.8},text=Muted,inner sep=0pt] at (46.5,80.3) {probability $p$};
  \node[font=\fsize{6.8},text=Muted,inner sep=0pt] at (75.1665,80.3) {probability $1-p$};
  \node[font=\fsize{7.4},inner sep=0pt] at (17.8335,83.6) {keep remainder};
  \node[font=\fsize{7.4},inner sep=0pt] at (46.5,83.6) {lose all};
  \node[font=\fsize{7.4},inner sep=0pt] at (75.1665,83.6) {keep money};

  % -------------------- (b) FOUR DESIGNS --------------------
  \path[rounded corners=1.8mm,fill=PanelFill] (94,1) rectangle (187,89);
  \begin{scope}
    \clip[rounded corners=1.8mm] (94,1) rectangle (187,89);
    \fill[HeadB] (94,1) rectangle (187,8.8);
  \end{scope}
  \ifheadrule \draw[HeadRuleB,line width=.55pt] (94,8.8) -- (187,8.8); \fi
  \path[rounded corners=1.8mm,draw=GreenLine,line width=.55pt]
    (94,1) rectangle (187,89);
  \node[anchor=west,text=HeadText,font=\fsize{14}\bfseries,inner sep=0pt]
    at (96.5,4.9) {\textcolor{HeadTagB}{(b)}\quad Four designs};

  % Model names are transcribed from the supplied image, not independently
  % verified release names. Logos are extracted symbols without raster labels.
  % Order follows the paper (19-09-2026): Haiku, Flash-Lite, Luna, Qwen, Grok.
  \card{96}{10.6}{185}{28.6}{ModelWash}{CardLine}
  \node[anchor=west,font=\fsize{9.0}\bfseries,inner sep=0pt]
    at (97.4,13.1) {Models we study};
  \card{97.4}{15.6}{113.7}{27.6}{ClaudeWash}{ClaudeLine}
  \card{115.0}{15.6}{131.3}{27.6}{GeminiWash}{GeminiLine}
  \card{132.6}{15.6}{148.9}{27.6}{GPTWash}{GPTLine}
  \card{150.2}{15.6}{166.5}{27.6}{QwenWash}{QwenLine}
  \card{167.8}{15.6}{184.1}{27.6}{GrokWash}{GrokLine}
  \asset{logo_claude}{105.55}{18.4}{5.6}{4.8}
  \asset{logo_gemini}{123.15}{18.4}{5.6}{4.8}
  \asset{logo_gpt}{140.75}{18.4}{5.6}{4.8}
  \asset{logo_qwen}{158.35}{18.4}{5.6}{4.8}
  \asset{logo_grok}{175.95}{18.4}{5.6}{4.8}
  \node[font=\fsize{7.3}\bfseries,inner sep=0pt] at (105.55,22.5) {Claude};
  \node[font=\fsize{7.3}\bfseries,inner sep=0pt] at (105.55,25.3) {Haiku 4.5};
  \node[font=\fsize{7.1}\bfseries,inner sep=0pt] at (123.15,22.5) {Gemini 3.5};
  \node[font=\fsize{7.1}\bfseries,inner sep=0pt] at (123.15,25.3) {Flash-Lite};
  \node[font=\fsize{7.3}\bfseries,inner sep=0pt] at (140.75,22.5) {GPT-5.6};
  \node[font=\fsize{7.3}\bfseries,inner sep=0pt] at (140.75,25.3) {Luna};
  \node[font=\fsize{7.0}\bfseries,inner sep=0pt] at (158.35,23.9) {Qwen3-235B};
  \node[font=\fsize{7.3}\bfseries,inner sep=0pt] at (175.95,23.9) {Grok 4.20};

  % Design cards use consistent geometry, title baselines, and artwork height.
  \card{96}{30.2}{139.5}{57.8}{SelfWash}{SelfLine}
  \card{141.5}{30.2}{185}{57.8}{PromptWash}{PromptLine}
  \card{96}{59.4}{139.5}{87}{ScriptWash}{ScriptLine}
  \card{141.5}{59.4}{185}{87}{MixedWash}{MixedLine}

  \designhead{101.4}{35.2}{1}{NumOne}{121.6}{Self-play}{same model $\times 6$}{NumTextOne}
  \designhead{146.9}{35.2}{2}{NumTwo}{167.1}{Prompt tests}{prompt variants}{NumTextTwo}
  \designhead{101.4}{64.4}{3}{NumThree}{121.6}{Scripted partners}{fixed partners}{NumTextThree}
  \designhead{146.9}{64.4}{4}{NumFour}{167.1}{Mixed tables}{two models}{NumTextFour}
  % All four groups have the same width (19.1 mm) so the robots are the same
  % size in every card, and each starts 8.8 mm below its card's top edge
  % (centre = top + 8.8 + half its height). The mixed-table PNG has no
  % transparent margin above its top antenna, so it sits 0.5 mm lower to
  % clear 'two models'. The scripted card keeps room under its robots for
  % the fifth Fixed label.
  \asset{group_self_play}{117.77}{46.56}{19.1}{19.1}
  \asset{group_prompt_tests}{163.25}{46.75}{19.1}{19.1}
  \asset{group_scripted_partners}{117.76}{76.48}{19.1}{19.1}
  \asset{group_mixed_tables}{163.25}{76.21}{19.1}{19.1}

  % Five Fixed labels are vector text, separated from the extracted artwork.
  % Four sit beside the upper and lower side robots, level with their faces;
  % the fifth is centred under the bottom robot, so the set is left-right
  % symmetric about the card's centre line (x = 117.76).
  \fixedbadge{105.5}{73.4}
  \fixedbadge{130.0}{73.4}
  \fixedbadge{105.5}{78.75}
  \fixedbadge{130.0}{78.75}
  \fixedbadge{117.76}{85.2}
\end{tikzpicture}
\end{document}
